% =============================================================================
%  Spectral Diagnostic Trajectories (SDT) - Standalone Abstract
% =============================================================================
\documentclass[10pt]{article}

\usepackage[T1]{fontenc}
\usepackage[utf8]{inputenc}
\usepackage{lmodern}
\usepackage[margin=0.85in, top=0.9in, bottom=0.9in]{geometry}
\usepackage{amsmath, amssymb}
\usepackage{xcolor}
\usepackage{microtype}
\usepackage{hyperref}

\hypersetup{
  colorlinks=true,
  linkcolor=blue!70!black,
  citecolor=blue!70!black,
  urlcolor=blue!70!black
}

\begin{document}
\thispagestyle{empty}

\begin{center}
  {\Large \textbf{Spectral Diagnostic Trajectories (SDT):}}\\[4pt]
  {\large \textbf{A Hypergraph-Spectral Framework for Adaptive Clinical}}\\[2pt]
  {\large \textbf{Decision Support and Operational Optimization in EHR Systems}}\\[12pt]
  {\textbf{Hemanth Manchabale Papachappa}}\thanks{Corresponding author: \texttt{hemanthmpkumar@gmail.com}} \quad
  {\textbf{Aniruddha Ganguly}}\\[4pt]
  {\small Independent Research in Clinical Informatics and Applied Mathematics}
\end{center}

\vspace{0.8em}

\begin{abstract}
\noindent
Emergency and critical care physicians must synthesize extensive, multi-modal electronic health record (EHR) data under stringent time constraints. While modern hospital systems capture dense longitudinal histories, structured laboratory orders, and continuous physiological telemetry, accessing decision-critical evidence during active patient workups remains cognitively demanding. Standard clinical information retrieval pipelines treat clinician interactions as isolated, stateless queries, overlooking the iterative, non-monotonic progression of diagnostic reasoning. When an initial working diagnosis is contradicted by new clinical findings---such as when an apparent community-acquired pneumonia abruptly deteriorates into an acute pulmonary embolism---conventional interfaces continue to surface results biased toward the earlier working hypothesis, contributing to cognitive anchoring, search distraction, and diagnostic delay.

\medskip
\noindent
Recent attempts to model clinical workflow continuity, such as the Geodesic Diagnostic Trajectories (GDT) framework, represent clinical intent as a continuous trajectory on the Riemannian manifold of Symmetric Positive Definite (SPD) covariance matrices. Although geometrically elegant, GDT incurs significant computational overhead ($O(d^3)$ per Riemannian matrix logarithm and exponential map), cannot natively represent multi-way interactions among disparate clinical modalities, and suffers from numerical instability when digesting high-frequency (125\,Hz) physiological telemetry streams.

\medskip
\noindent
In this work, we present \textbf{Spectral Diagnostic Trajectories (SDT)}, a graph-theoretic framework that models the diagnostic search space as an evolving \emph{multi-modal hypergraph}. In SDT, clinical text, structured events, and continuous waveforms (encoded via Mamba selective state-space models) are represented as heterogeneous nodes interconnected by multi-way hyperedges. Our formulation introduces four methodological components:
(1)~\emph{Spectral Pivot Detection}, which tracks the algebraic connectivity (Fiedler value $\lambda_2$) of the normalized hypergraph Laplacian in real time, identifying hypothesis shifts in $O(|E_t|)$ time;
(2)~\emph{Graph-Signal Context Attenuation}, which employs a Chebyshev high-pass filter and the Fiedler vector to sever obsolete diagnostic subgraphs with provable Cheeger conductance bounds;
(3)~\emph{Spectrally-Biased Anticipatory Prefetching}, formulating prospective record retrieval as a biased random walk over the Laplacian eigenspace of the institutional patient graph; and
(4)~\emph{Cognitive-Load-Bounded Delivery}, which selects prefetched records via a greedy 0/1 knapsack optimization based on Sweller's cognitive load theory.

\medskip
\noindent
In a nine-arm within-condition repeated-measures simulation across 1{,}000 standardized diagnostic vignettes (9{,}000 evaluated sessions) calibrated against empirical EHR workflow studies, SDT reduced mean time-to-correct-information (TTCI) by 28.9\% relative to standard search (62.97\,s vs.\ 88.52\,s) and by 8.7\% relative to GDT (68.96\,s), while improving diagnostic retrieval accuracy from 86.8\% to 93.9\% and decreasing modeled cognitive load by 32.2\%. Average algorithmic response latency was 112.9\,ms. Translating these findings into an Erlang~C ($M/M/c$) emergency department operational model demonstrates that individual per-encounter time savings compound non-linearly, yielding an estimated 12.3-minute reduction in patient queue waiting times under near-saturation conditions.
\end{abstract}

\vspace{1em}
\noindent\textbf{Keywords:}
Hypergraph Neural Networks; Spectral Graph Theory; Algebraic Connectivity;
Fiedler Vector; Clinical Decision Support; EHR Information Retrieval;
Cognitive Load Theory; NASA-TLX; Erlang C Queueing; Mamba State-Space Models.

\end{document}
